\subsection{截面的喷射}

12.5.1. 设 \(\pi: Y \to X\) 是浸没，令 \(x \in X\)，并令 \(s \in J_x^k(X, Y)\)。我们称 \(s\) 为 \(k\) 阶的 \(\pi\) 的截面喷射，即 \(s\) 具有 \(j_x^k(f)\) 的形式，其中 \(f\) 是类 \(C^r\) 的 \(\pi\) 的截面，定义在 \(x\) 的某个开邻域上；这一条件等价于

\[
j _ {b (s)} ^ {k} (\pi) \circ s = j _ {x} ^ {k} (\mathrm{Id} _ {\mathrm{x}}).
\]

我们把 \(\mathbf{P}_x^k (\pi)\)（相应地 \(\mathbf{P}^k (\pi)\)）记作由截面喷射组成的集合，其中 \(s\in \mathrm{J}_x^k (\mathrm{X},\mathrm{Y})\)（相应地 \(\mathrm{J}^k (\mathrm{X},\mathrm{Y}))\) 的元素是 \(\pi\) 的截面喷射；它是类 \(\mathbf{C}^{r - k}\) 的 \(\mathrm{J}^k (\mathrm{X},\mathrm{Y})\) 子流形。映射

\[
s \colon \mathrm{P} ^ {k} (\pi) \rightarrow \mathrm{X} \quad \text {and} \quad b \colon \mathrm{P} ^ {k} (\pi) \rightarrow \mathrm{Y}
\]

是浸没；若 \(\pi\) 是纤维化，则它们是纤维化。当 \(k=0\) 时，\(b\) 是同构，据此将 \(\mathrm{P}^{0}(\pi)\) 与 Y 等同。

若 \(k' \leqslant k\)，映射 \(r^{k, k'} : \mathrm{J}^k(\mathrm{X}, \mathrm{Y}) \to \mathrm{J}^{k'}(\mathrm{X}, \mathrm{Y})\) 将 \(\mathrm{P}^k(\pi)\) 映到 \(\mathrm{P}^{k'}(\pi)\)；由此得到的从 \(\mathrm{P}^k(\pi)\) 到 \(\mathrm{P}^{k'}(\pi)\) 的映射 \(r^{k, k'}\) 是类 \(\mathbf{C}^{r-k}\) 的纤维化。

12.5.2. 设 Z 是类 \(C^{r}\) 的流形；假设 \(Y = X \times Z\)，且 \(\pi = pr_{1}: Y \to X\)。将 \(J^{k}(Id_{X}, pr_{2})\)（参见 12.3.9）限制在 \(P^{k}(\pi)\) 上所得的映射，是从 \(P^{k}(\pi)\) 到 \(J^{k}(X, Z)\) 的同构，借此将这两个流形等同。

12.5.3. 设 \(\mathbf{Y}'\) 是类 \(\mathbf{C}^{r}\) 的流形，令 \(\pi: \mathrm{Y} \to \mathrm{X}\) 和 \(\pi': \mathrm{Y}' \to \mathrm{X}\) 为浸没，且令 \(g: \mathrm{Y} \to \mathrm{Y}'\) 为满足 \(\pi' \circ g = \pi\) 的态射。映射 \(\mathbf{J}^{k}(\mathrm{Id}_{\mathbf{X}},g)\) 的限制得到一个映射

\[
\mathbf {P} ^ {k} (g) \colon \mathbf {P} ^ {k} (\pi) \to \mathbf {P} ^ {k} (\pi^ {\prime})
\]

它属于类 \(C^{r-k}\)。

设 \(\mathbf{X}'\) 是类 \(\mathbf{C}^r\) 的流形，且 \(f: \mathbf{X}' \to \mathbf{X}\) 是态射；置 \((\mathbf{Y}', \pi') = f^*(\mathbf{Y}, \pi)\)（参见 5.11.5）。映射 \(\mathbf{J}^k(f, \mathrm{Id}_\mathbf{Y})\) 定义了一个类 \(\mathbf{C}^{r-k}\) 的映射，从 \(f^*\mathbf{P}^k(\pi)\) 到 \(\mathbf{P}^k(\pi')\)。

